\documentclass{article}
\usepackage{pathfinder-readable}

\title{When a Fixed-Batch Risk Estimate Changes an Obedience Test}

\begin{document}
\maketitle

\begin{abstract}
The supplied pair consists of a paper on rewards for AI agents and a local paper on fixed-batch risk estimates
in a resource-allocation game. We asked whether an estimated risk value can change the first paper's test for
supporting a signal-dependent action policy with rewards that do not depend on the state. For every finite
batch size, the peers found a two-signal example that passes the test under the expected empirical estimate
and fails it under true conditional risk. They also checked a margin test and a common-reward test for bounded
value errors. The verifier accepted these narrow results but left the thread in DRAFT because the report stage
under a specified risk timing, realised batches, and a broader novelty claim remain open.
\end{abstract}

\section{Introduction}

The supplied Q, by Bergemann, Koh and Morris, studies how to reward an AI agent when its preferences and
abilities are private. The agent may report a type and then choose an action after receiving a signal about
the world. Q asks when rewards can make both the report and the action follow a desired policy. In the special
case where an action's reward cannot depend on the state, Q gives a test based on cycles of signals: the gains
from following the policy around every such cycle must add to at least zero \cite{Q}.

The supplied P studies a different setting: a quadratic resource-allocation tax game evaluated with
conditional value-at-risk, or CVaR. CVaR measures the average loss in a specified upper tail. P separates true
CVaR from the average empirical CVaR produced by repeatedly drawing a fixed-size batch. Its one-sample example
has an exact equilibrium for the empirical objective that fails participation under true risk. Its value-error
bounds concern the expected empirical objective, subject to its bounded-loss and valid-query assumptions
\cite{P}.

P calls its resource-allocation source ``Q'', but that is a different paper from the Q supplied here. The
peers therefore pursued a new, narrower connection: could P's sampling bias change Q's signal-cycle answer?
They treated negative \emph{conditional} CVaR, evaluated after a signal, as the agent's action utility in Q.
This is an explicit modelling choice, not a result about the tax game in P or about CVaR of an entire plan
before the signal.

\section{What was done}

The peers first checked the identity of the papers and set aside a direct transfer of P's tax results. They
then reduced Q to one agent type, two equally likely states that the signal reveals perfectly, and two
actions, \(a\) and \(b\). They let a fresh batch estimate the risk of each action after the signal. Ada
proposed the reversed-risk example; Emmy checked it against Q's cycle condition and extended it from one
observation to every finite batch size. Both peers checked the reward inequalities independently. Ada then
derived a test for a reward that works across a box of possible utility errors, and Emmy checked it. Emmy
separately tested what happens if the agent instead assesses a whole plan by CVaR before seeing the signal.

For the main example, fix a positive loss scale \(h\) and a finite batch size \(m\). A risky action loses
either \(0\) or \(2h\), each with probability one half. A safe action has constant loss \(c_m\). In state
\(0\), \(a\) is risky and \(b\) is safe; in state \(1\), their roles reverse. The proposed policy selects the
risky action in both states. Let \(q_m\) be the expected empirical upper-half CVaR of the risky loss from
\(m\) independent observations. An all-zero batch has probability \(2^{-m}\), so its empirical CVaR is zero;
every other empirical value is at most \(2h\). The peers chose
\[
q_m\leq 2h(1-2^{-m})<c_m:=2h-h2^{-m}<2h.
\]
The true upper-half CVaR of the risky action is \(2h\). These inequalities put the safe loss between the
estimated and true values of the risky loss.

\section{Results}

Under expected empirical CVaR, the proposed risky action has lower estimated loss than the safe one after
either signal. A reward of zero for each action therefore supports the policy. Under true conditional CVaR,
let \(r(a)\) and \(r(b)\) be rewards that depend only on the action. Obedience in state \(0\) requires
\(r(a)-r(b)\geq 2h-c_m\); obedience in state \(1\) requires \(r(b)-r(a)\geq 2h-c_m\). Since
\(2h-c_m=h2^{-m}>0\), both requirements cannot hold. Equivalently, Q's two-signal cycle has a positive
estimated sum and the true sum \(-2h2^{-m}\). Thus, for every finite \(m\), some instance passes the expected
empirical test but no state-independent reward supports the same policy under true conditional risk. The safe
loss \(c_m\) varies with \(m\); this is not failure for one fixed instance as the batch grows. Ada's
derivation in \texttt{ada/finite\_batch\_implementability.md} and Emmy's check in \texttt{emmy/findings.md}
support this result.

The peers also checked when value accuracy protects the cycle test. Write \(U(x,s)\) for true utility of
action \(x\) after signal \(s\), and \(\widetilde U(x,s)\) for its estimated utility. If \(|U(x,s)-\widetilde
U(x,s)|\leq e\) for every compared action and signal, each term of a signal cycle can change by at most
\(2e\). Hence an estimated cycle of length \(K\) with a margin of at least \(2Ke\) still meets Q's true cycle
condition. If all relevant simple cycles have that margin, Q's result supplies some supporting
state-independent reward; the estimated reward itself need not work under true utility.

A stronger question asks for \emph{one} reward that works for every utility array in an independent error box
of width \(e\). Let \(B\) be the actions prescribed somewhere by the policy, and let \(a(s)\) be its action
after signal \(s\). For distinct \(x,y\in B\), the peers assigned the directed edge \(x\) to \(y\) the weight
\[
w(x,y)=\max_{s:\,a(s)=x}\bigl\{\widetilde U(y,s)-\widetilde U(x,s)+2e\bigr\}.
\]
Their check showed that a common state-independent reward exists exactly when no directed cycle of these
prescribed actions has positive total weight. The reason is that robust obedience says \(r(x)-r(y)\geq
w(x,y)\). Summing those inequalities rules out a positive cycle, while Q's path construction gives rewards
when none exists. Actions never prescribed can be prohibited. This exact statement concerns the independent
error box; using P's bound for a particular application still requires P's assumptions.

\section{Objections and limits}

The main example uses risk assessed after each signal. Emmy checked a separate timing objection: an agent may
instead rank a whole signal-dependent plan by CVaR of the joint signal-and-loss distribution before seeing a
signal. In a two-signal example in the ledger, averaging conditional CVaRs favoured one report, while ex-ante
CVaR favoured another. Thus, averaging local risk values does not generally give Q's report-stage continuation
value. More generally, if a loss \(J\) lies between \(0\) and \(M\), \(S\) is the signal, and \(\beta\) is
upper-tail mass, the peers established
\[
0\leq \operatorname{CVaR}_{\beta}(J)
-\mathbb E\!\left[\operatorname{CVaR}_{\beta}(J\mid S)\right]
\leq (1-\beta)M.
\]
This difference can remain with exact risk values. The ledger cites Miller and Yang \cite{MillerYang} and
Rudloff, Street and Vallad\~ao \cite{RudloffStreetValladao} on the known dynamic-CVaR timing issue. Its
limited searches did not settle novelty.

Q also allows state-dependent rewards outside the subclass studied here. The example does not rule those out.
It concerns the \emph{expected} empirical value, not a guarantee for one realised batch or convergence of
strategic learning. The peers established a two-signal obstruction and value-error tests, not a full theorem
combining Q's action and type-report cycles under CVaR.

\section{Why the thread stopped}

The verifier's single round accepted the finite-batch reversal, the reward inequalities, and the value-error
graph test. It recorded DRAFT because the supported conclusion is narrow: the peers had not fixed a
risk-timing model for the report stage or proved the corresponding nested-cycle result. The smallest next step
is to specify whether reports are evaluated by ex-ante plan CVaR or by conditional CVaR after each signal,
then solve a two-type, two-signal report game under that choice. A realised-batch guarantee and a fuller
prior-work check would remain separate work.

\bibliographystyle{plainurl}
\bibliography{references}
\end{document}
